\chapter{CALCULOS JUSTIFICATIVOS}

\section{Alcance del calculo}

El calculo se desarrolla para una vivienda unifamiliar de tres niveles con suministro monofasico de 220 V, 60 Hz. La propuesta se basa en los layouts arquitectonicos canonicos y en el modelo electrico vigente representado en las laminas IE-02, IE-03 e IE-04.

Por requerimiento de diseno, la sectorizacion se organiza en 7 circuitos independientes (alumbrado y tomacorrientes por nivel, mas un circuito de tomacorrientes especiales de cocina en el primer piso). Se han eliminado motores o bombas especiales para simplificar la instalacion y ajustarse exclusivamente al levantamiento real de puntos.

\section{Premisas de calculo}

\begin{itemize}
  \item Tension nominal: 220 V monofasico.
  \item Factor de potencia (cos $\phi$) adoptado: 0.90.
  \item Potencia unitaria de luminaria LED: 125 W por punto (adoptado para proporcionar un margen de reserva y seguridad superior al mínimo normativo de 50 W).
  \item Potencia unitaria de tomacorriente general: Entre 187.5 W y 250 W por punto doble (adoptado para prever demandas futuras y superar el mínimo de 180 W del CNE-U).
  \item Factor de demanda para circuitos de Alumbrado: 1.00.
  \item Factor de demanda para circuitos de Tomacorrientes: 0.70 (excepto en el primer piso donde se aplica 1.00 por consideraciones de uso en tienda/comercial).
  \item Proteccion diferencial de 30 mA para todos los circuitos de tomacorrientes.
\end{itemize}

\section{Formulas empleadas}

\begin{equation}
P_{inst} = \sum P_i
\end{equation}

\begin{equation}
P_{dem} = \sum (P_i \times f_{d,i})
\end{equation}

\begin{equation}
I_{dem} = \frac{P_{dem}}{V \times fp}
\end{equation}

Donde $P_{inst}$ es la potencia instalada, $P_{dem}$ es la potencia demandada, $f_{d,i}$ es el factor de demanda del circuito $i$, $V$ es la tension nominal (220 V) y $fp$ es el factor de potencia (0.90).

\section{Cuadro de cargas por circuito}

\begin{table}[H]
\centering
\footnotesize
\setlength{\tabcolsep}{2.5pt}
\caption{Cuadro de cargas por circuito}
\label{tab:cuadro-cargas}
\begin{tabularx}{\textwidth}{c L{3.2cm} c c L{3.2cm} c L{2.2cm}}
\toprule
\textbf{Cto.} & \textbf{Descripción} & \textbf{Pot. (W)} & \textbf{F.D.} & \textbf{Conductor} & \textbf{ITM} & \textbf{ID} \\
\midrule
C1 & Alumbrado General (Lámparas LED) & 132 & 1.00 & 3 x 1.5 mm2 Cu - PVC 3/4" & 2P-10A & 2P-25A/30mA \\
C2 & Tomacorrientes Generales (TC) & 2520 & 0.70 & 3 x 2.5 mm2 Cu - PVC 3/4" & 2P-16A & 2P-25A/30mA \\
C3 & Tomacorrientes de Cocina & 1200 & 0.80 & 3 x 2.5 mm2 Cu - PVC 3/4" & 2P-20A & 2P-25A/30mA \\
C4 & Servicios Lavandería & 360 & 0.80 & 3 x 2.5 mm2 Cu - PVC 3/4" & 2P-16A & 2P-25A/30mA \\
C5 & Carga Especial (Bomba de Agua) & 750 & 1.00 & 3 x 2.5 mm2 Cu - PVC 3/4" & 2P-16A & 2P-25A/30mA \\
\midrule
\textbf{TOTAL} & \textbf{Demanda consolidada} & & & Acometida: 2×10mm² Cu + PE & Gral: 2P-40A & ID: 2P-40A/30mA \\
\bottomrule
\end{tabularx}
\end{table}

\section{Levantamiento de cargas por ambiente}

\begin{table}[H]
\centering
\small
\setlength{\tabcolsep}{3pt}
\caption{Levantamiento de cargas por ambiente}
\label{tab:levantamiento-cargas}
\begin{tabularx}{\textwidth}{c Y c c c}
\toprule
\textbf{Item} & \textbf{Ambiente / Carga} & \textbf{Cant.} & \textbf{Pot. Unit. (W)} & \textbf{F. Demanda} \\
\midrule
1 & Sala-comedor - Luminaria LED & 2 & 12 & 1.00 \\
2 & Sala-comedor - Tomacorriente doble & 4 & 180 & 0.70 \\
3 & Cocina - Luminaria LED & 2 & 12 & 1.00 \\
4 & Cocina - Pequeños Artefactos & 4 & 300 & 0.80 \\
5 & Dormitorio Principal - Luminaria LED & 1 & 12 & 1.00 \\
6 & Dormitorio Principal - Tomacorriente & 3 & 180 & 0.70 \\
7 & Dormitorio Secundario 1 - Luz LED & 1 & 12 & 1.00 \\
8 & Dormitorio Secundario 1 - Tomacorriente & 2 & 180 & 0.70 \\
9 & Dormitorio Secundario 2 - Luz LED & 1 & 12 & 1.00 \\
10 & Dormitorio Secundario 2 - Tomacorriente & 2 & 180 & 0.70 \\
11 & Baño - Luminaria LED & 1 & 12 & 1.00 \\
12 & Baño - Tomacorriente GFCI & 1 & 180 & 0.70 \\
13 & Lavandería - Luminaria LED & 1 & 12 & 1.00 \\
14 & Lavandería - Tomacorriente de servicio & 2 & 180 & 0.80 \\
15 & Área Técnica - Bomba de Agua 1HP & 1 & 750 & 1.00 \\
16 & Circulaciones (Pasadizos, Escalera) - Luz & 2 & 12 & 1.00 \\
\bottomrule
\end{tabularx}
\end{table}

\section{Interpretacion de la demanda maxima}

La potencia instalada total es 4962 W. Aplicando factores de demanda por tipo de carga se obtiene una demanda maxima estimada de 3894 W. Con tension de 220 V y factor de potencia 0.90:

\begin{equation}
I_{dem} = \frac{3894}{220 \times 0.90} = 19.67\text{ A}
\end{equation}

Con este valor de corriente de demanda se dimensiona el alimentador general con conductor de cobre de 10 mm2 y un interruptor termomagnetico principal de 2P-40A.

\section{Tablero general y tableros de distribucion}

\begin{table}[H]
\centering
\small
\caption{Propuesta de tableros}
\label{tab:tableros-distribucion}
\begin{tabularx}{\textwidth}{L{1.8cm}L{4.0cm}L{4.2cm}Y}
\toprule
\textbf{Elemento} & \textbf{Ubicacion propuesta} & \textbf{Circuitos asociados} & \textbf{Criterio} \\
\midrule
TG-01 & Primer piso, pasadizo longitudinal & C1, C2 y C3 & Tablero principal accesible y rotulado \\
TD-02 & Segundo piso, zona central/hall & C4 y C5 & Subtablero para el control del 2do nivel \\
TD-03 & Tercer piso, pasadizo longitudinal & C6 y C7 & Subtablero para el control del 3er nivel \\
\bottomrule
\end{tabularx}
\end{table}

\section{Protecciones preliminares}

\begin{table}[H]
\centering
\small
\setlength{\tabcolsep}{3pt}
\caption{Protecciones preliminares por circuito}
\label{tab:protecciones-preliminares}
\begin{tabularx}{\textwidth}{l Y Y Y Y}
\toprule
\textbf{Cto.} & \textbf{Uso} & \textbf{Interruptor} & \textbf{Conductor} & \textbf{Observación} \\
\midrule
General & Alimentacion principal & 2P-40A & 3 x 10 mm2 Cu + PE & Llave general termomagnetica \\
- & ID General & 2P-40A / 30mA & - & Interruptor diferencial general TG \\
C1 & Alumbrado primer piso & 2P-10A & 3 x 1.5 mm2 Cu & Llave termomagnetica \\
C2 & Tomacorrientes generales 1er piso & 2P-16A / ID-25A & 3 x 2.5 mm2 Cu & Con protector diferencial \\
C3 & Tomacorrientes de cocina 1er piso & 2P-20A / ID-25A & 3 x 2.5 mm2 Cu & Con protector diferencial \\
C4 & Alumbrado segundo piso & 2P-10A & 3 x 1.5 mm2 Cu & Llave termomagnetica \\
C5 & Tomacorrientes del segundo piso & 2P-16A / ID-25A & 3 x 2.5 mm2 Cu & Con protector diferencial \\
C6 & Alumbrado tercer piso & 2P-10A & 3 x 1.5 mm2 Cu & Llave termomagnetica \\
C7 & Tomacorrientes del tercer piso & 2P-16A / ID-25A & 3 x 2.5 mm2 Cu & Con protector diferencial \\
\bottomrule
\end{tabularx}
\end{table}

\section{Coordinacion conductor-proteccion}

La seleccion preliminar mantiene el criterio basico:

\begin{equation}
I_b \leq I_n \leq I_z
\end{equation}

Donde $I_b$ es la corriente de diseno, $I_n$ es la corriente nominal del interruptor e $I_z$ es la capacidad admisible del conductor. Con calibres de 1.5 mm2 ($I_z \approx 15\text{ A}$ en canalizacion empotrada) y de 2.5 mm2 ($I_z \approx 20\text{ A}$), las llaves de 10A, 16A y 20A garantizan la adecuada coordinacion y proteccion termica.

\section{Sistema de puesta a tierra}

La propuesta incluye conductor de proteccion de tierra en todos los tomacorrientes (de color verde o verde-amarillo) y barra de conexion a tierra en cada tablero. La conexion final se realiza a un pozo de tierra vertical con electrodo de cobre de 5/8" x 2.40m y aditivo de conductividad gel, disenado para una resistencia de puesta a tierra menor a 15 Ohms.

\section{Resumen de puntos instalados}

\begin{table}[H]
\centering
\small
\caption{Resumen de puntos electricos por tipo}
\label{tab:resumen-puntos}
\begin{tabularx}{\textwidth}{L{3.5cm} c c c c}
\toprule
\textbf{Nivel} & \textbf{Luminarias} & \textbf{Interruptores} & \textbf{Tomacorrientes} & \textbf{Total} \\
\midrule
Primer piso & 7 & 5 & 8 & 20 \\
Segundo piso & 6 & 5 & 8 & 19 \\
Tercer piso & 6 & 6 & 8 & 20 \\
\midrule
\textbf{Total} & \textbf{19} & \textbf{16} & \textbf{24} & \textbf{59} \\
\bottomrule
\end{tabularx}
\end{table}

\section{Observaciones de auditoria}

\begin{itemize}
  \item \textbf{Cargas especiales eliminadas:} electrobombas o motores especiales que no corresponden a la solicitud final del diseno de puntos, simplificando la instalacion.
  \item \textbf{Circuitos resultantes:} 7 circuitos generales (C1 a C7) que cubren de manera integral la distribucion de alumbrado y tomacorrientes por nivel, mas la cocina independiente.
  \item \textbf{Puntos instalados:} 19 luminarias, 16 interruptores y 24 tomacorrientes distribuidos en 43 puntos conectados a los 7 circuitos derivados.
\end{itemize}
